\documentclass[11pt]{article}
\usepackage[a4paper,margin=24mm]{geometry}
\usepackage{amsmath,amssymb,amsthm,mathtools}
\usepackage[hidelinks]{hyperref}
\newtheorem{theorem}{Theorem}
\newtheorem{lemma}{Lemma}
\newcommand{\CF}{\mathop{\mathrm K}}
\title{A proof of the four-pair gamma-ratio\\continued-fraction conjecture\\MathDB 352751}
\author{Principia Math In Progress}
\date{18 September 2026}
\begin{document}
\maketitle
\begin{abstract}
We prove Conjecture 9.3 of Cao--Tanigawa--Zhai. A single
Bauer--Muir transformation, with a constant adjoint factor, gives
a shift by two combined with a sign change of one parameter.
The resulting gamma identity follows from a periodic quotient
and its limit at infinity. An elementary tail convergence argument
justifies the entire stated complex half-plane; termination handles
the remaining cases.
\end{abstract}

\section*{Task 1: statement and conventions}
Write
\begin{align}
P(x,l,n,\eta)
={}&
\frac{\Gamma(\frac{x+l+n+\eta+1}{4})
      \Gamma(\frac{x+l-n-\eta+1}{4})}
     {\Gamma(\frac{x+l+n-\eta+3}{4})
      \Gamma(\frac{x+l-n+\eta+3}{4})}
\notag\\[-2pt]
&{}\times
\frac{\Gamma(\frac{x-l+n-\eta+3}{4})
      \Gamma(\frac{x-l-n+\eta+3}{4})}
     {\Gamma(\frac{x-l+n+\eta+1}{4})
      \Gamma(\frac{x-l-n-\eta+1}{4})}.
\label{P}
\end{align}
For \(m\ge1\), set
\begin{align}
a_{2m}&=(2m)^2-l^2,&
a_{2m-1}&=\frac{((2m-1)^2-n^2)((2m-1)^2-\eta^2)}
                    {(2m-1)^2},\label{a}\\
b_{2m}&=x-\frac{n\eta}{2m+1},&
b_{2m-1}&=x+\frac{n\eta}{2m-1},\qquad b_0=x-n\eta.
\label{b}
\end{align}
Define \(F=b_0+\CF_{k=1}^{\infty}(a_k/b_k)\), using ordinary
finite convergents. A zero partial numerator terminates the fraction.
As usual for gamma identities with complex parameters, the displayed
formulas are meromorphic identities and give ordinary values wherever
the expressions are defined. Poles and removable singularities are
understood by continuation.

\begin{theorem}\label{main}
For arbitrary \(l,n,\eta\in\mathbb C\) and \(\Re x>0\),
\begin{equation}
 \frac{1-P(x,l,n,\eta)}{1+P(x,l,n,\eta)}
     =\frac{l}{F(x,l,n,\eta)}.
\label{answer}
\end{equation}
The nonterminating continued fraction converges in the extended
complex plane, locally uniformly away from its poles.
If \(n\) or \(\eta\) is an odd integer, or \(l\) is an even integer,
the identity extends to arbitrary complex \(x\) with the same
meromorphic convention. This proves the full domain asserted in
\cite[Conjecture 9.3, p.~30]{CTZ}.
\end{theorem}

\subsection*{An elementary convergence lemma}
The following lemma supplies the analytic justification, rather
than treating continuation of the answer as convergence of its
continued fraction.

\begin{lemma}\label{convergence}
Suppose the coefficients of a continued fraction depend
holomorphically on finitely many parameters and, uniformly on compact
parameter sets, satisfy
\[
 a_k=k^2+O(1),\qquad b_k=x+O(k^{-1}),\qquad \Re x>0.
\]
Where no partial numerator vanishes, its convergents have a
meromorphic limit, with ordinary convergence in the extended
complex plane and local uniform convergence away from poles.
\end{lemma}
\begin{proof}
Work on a relatively compact parameter polydisc with \(\Re x>0\).
For sufficiently large \(k\), the analytic logarithm
\[
 L_k=\log(a_k/k^2)=O(k^{-2})
\]
is unambiguously the branch near zero. Define
\[
 t_k=\sum_{j=1}^{\infty}(-1)^j L_{k+j},\qquad s_k=e^{t_k}.
\]
This series converges absolutely and locally uniformly,
\(t_k=O(k^{-1})\), and \(t_{k-1}+t_k=-L_k\).
The equivalence transformation
\[
 \widetilde a_k=s_{k-1}s_ka_k=k^2,\qquad
 \beta_k=\widetilde b_k=s_kb_k=x+O(k^{-1})
\]
therefore gives \(\Re\beta_k\ge\delta>0\) for all \(k\ge K\),
after increasing \(K\). The transformation multiplies the tail
starting at \(b_K\) by \(s_K\) and is exact for every finite truncation.

Consider that transformed tail through \(N\). Let \(T_N\) be its
tridiagonal matrix on indices \(K,\ldots,N\), with diagonal
\(\beta_k\) and entries \(i(k+1)\) in both positions connecting
\(k\) to \(k+1\). Its off-diagonal part is skew-Hermitian. Thus
\[
 \Re\langle T_Nv,v\rangle
       =\sum_{k=K}^N\Re\beta_k\,|v_k|^2
       \ge\delta\|v\|_2^2.
\]
In particular \(T_N\) is invertible and the vector
\(v^{(N)}=T_N^{-1}e_K\), extended by zeros, satisfies
\(\|v^{(N)}\|_2\le\delta^{-1}\). Schur complementation gives
\[
 v_K^{(N)}
 =\frac{1}{\beta_K+
       \dfrac{(K+1)^2}{\beta_{K+1}+\cdots+
       \dfrac{N^2}{\beta_N}}}.
\tag{*}
\]
Every weakly convergent subsequence has an \(\ell^2\) limit solving
the infinite tridiagonal equation \(Tv=e_K\), since each coordinate
equation uses only three coordinates.

There is at most one such \(\ell^2\) solution. Indeed, if \(Tw=0\)
and \(w\ne0\), summing the real parts of its equations through \(N\)
gives
\[
 \delta\sum_{k=K}^N|w_k|^2
 \le (N+1)|w_Nw_{N+1}|.
\]
The left side is bounded below by a positive constant for every
sufficiently large \(N\). Hence
\[
 |w_N|^2+|w_{N+1}|^2\ge\frac{c}{N+1}
\]
for some \(c>0\), contradicting square summability. This also shows
that all weak subsequential limits of \(v^{(N)}\) coincide.
Consequently every coordinate, in particular \((*)\), converges.
The limiting vector is nonzero because it solves \(Tv=e_K\).
Moreover the finite energy identity implies
\[
 \Re v_K\ge\delta\|v\|_2^2>0,
\]
by taking limits of any fixed finite subsum. Thus the limiting
reciprocal tail is nonzero.

The functions \(v_K^{(N)}\) are holomorphic on the polydisc and
uniformly bounded by \(\delta^{-1}\). Normal-family compactness and
the pointwise uniqueness just proved yield locally uniform
holomorphic convergence. Undo the equivalence transformation and
attach the finite initial segment. This applies a finite
nondegenerate linear fractional transformation, and proves the
claimed meromorphic and extended-plane convergence.
\end{proof}

For \eqref{a}--\eqref{b}, both asymptotic hypotheses hold uniformly
on compact subsets of \(\{\Re x>0\}\times\mathbb C^3\).
If a partial numerator vanishes, the finite terminated expression
is meromorphic directly. The tail construction above also provides
meromorphic dependence near such parameter values after attaching
the finite prefix, with the usual removals wherever defined.

\subsection*{One Bauer--Muir transformation}
Initially take real parameters
\[
 x>3,\qquad 0<l<\tfrac14,\qquad
 |n|<\tfrac14,\quad |\eta|<\tfrac14.
\tag{D}
\]
All \(a_k,b_k,b_0\) are then positive. The lemma gives convergence,
and positivity gives the useful bound
\begin{equation}
 x-n\eta\le F(x,l,n,\eta)
      \le x-n\eta+\frac{(1-n^2)(1-\eta^2)}{x+n\eta}.
\label{bound}
\end{equation}
Put
\begin{align}
 u&=\frac{l^2-(n+\eta)^2-x^2+1}{2(x+1)},&
 c&=x+u=\frac{(x+1)^2+l^2-(n+\eta)^2}{2(x+1)},\label{uc}\\
 r_{2m}&=2m+u+\frac{n\eta}{2m+1}\quad(m\ge0),&
 r_{2m-1}&=2m-u-x\quad(m\ge1).\label{r}
\end{align}
The adjoint factors are all the same:
\begin{equation}
 a_k-r_{k-1}(b_k+r_k)=c^2-l^2=:D.
\label{adjoint}
\end{equation}
To check this without induction, for even \(k\) the left side is
\[
 k^2-l^2-(k-u-x)(k+u+x)=(u+x)^2-l^2.
\]
For odd \(k\) it is
\[
 \frac{(k^2-n^2)(k^2-\eta^2)}{k^2}
 -(k-1+u+n\eta/k)(k+1-u+n\eta/k)
 =(u-1)^2-(n+\eta)^2.
\]
These expressions agree by \eqref{uc}; \(D\ne0\) on (D).
The transformed denominators satisfy
\begin{align}
 b_0+r_0&=c,&
 b_1+r_1&=b_0(x+2,l,n,-\eta)-c,\notag\\
 b_k+r_k-r_{k-2}&=b_{k-1}(x+2,l,n,-\eta)\quad(k\ge2).
\label{denominators}
\end{align}
Since the ratios of consecutive adjoint factors equal one,
the classical Bauer--Muir identity
\cite[Lemma 3.1, p.~8]{CTZ} consequently gives
\begin{equation}
 F(x,l,n,\eta)
   =c+\frac{c^2-l^2}{F(x+2,l,n,-\eta)-c}.
\label{BM}
\end{equation}

For completeness, the limiting step in this use of the transformation
is justified directly. If \(A_N/B_N\) are the original convergents,
the finite transformed fraction is
\[
 \frac{A_N+r_NA_{N-1}}{B_N+r_NB_{N-1}}
 =c+\frac{D}{F_{N-1}(x+2,l,n,-\eta)-c}.
\]
This finite identity follows from the continuant recurrence and
\eqref{adjoint}. On (D), \(B_N>0\), and \(r_N>0\) for all sufficiently
large \(N\). Its left side is therefore a convex combination of two
successive convergents and tends to \(F\). The transformed tail
converges by Lemma~\ref{convergence}; this proves \eqref{BM}, or its
cross-multiplied version if necessary. No assumption about arbitrary
modified approximants sharing a limit is being made.

\subsection*{Identification of the gamma quotient}
Suppress \(l,n\) from the notation and set
\[
 Q(x,\eta)=\frac{F(x,l,n,\eta)-l}{F(x,l,n,\eta)+l}.
\]
Equation \eqref{BM} says
\[
 FF'-c(F+F')+l^2=0,\qquad F'=F(x+2,l,n,-\eta).
\]
Consequently
\begin{equation}
 Q(x,\eta)Q(x+2,-\eta)
 =\frac{c-l}{c+l}
 =\frac{(x+1-l)^2-(n+\eta)^2}
        {(x+1+l)^2-(n+\eta)^2}
 =:R(x,\eta).
\label{Qrec}
\end{equation}
All the quantities used in division are nonzero on (D).

The gamma quotient in \eqref{P} satisfies precisely the same identity:
\begin{equation}
 P(x,l,n,\eta)P(x+2,l,n,-\eta)=R(x,\eta).
\label{Prec}
\end{equation}
Indeed, multiply the eight gamma factors in the two numerators and
denominators. Four pairs cancel directly. The remaining four ratios use
\(\Gamma(A+1)/\Gamma(A)=A\), giving
\[
 \frac{(x+1-l+n+\eta)(x+1-l-n-\eta)}
      {(x+1+l+n+\eta)(x+1+l-n-\eta)}.
\]
The standard gamma-ratio asymptotic
\(\Gamma(y+a)/\Gamma(y+b)\sim y^{a-b}\), or Euler's gamma limit,
gives \(P(x,l,n,\eta)\to1\) as real \(x\to+\infty\):
the sums of the shifts in its numerator and denominator agree.
Also \eqref{bound} gives \(F=x-n\eta+O(x^{-1})\), hence \(Q\to1\).

Now \(H(x,\eta)=Q(x,\eta)/P(x,l,n,\eta)\) obeys
\[
 H(x,\eta)H(x+2,-\eta)=1,\qquad
 H(x+2,-\eta)H(x+4,\eta)=1.
\]
Thus \(H(x,\eta)=H(x+4,\eta)\). For any positive integer \(j\),
\[
 H(x,\eta)=H(x+4j,\eta)\longrightarrow1,
\]
so \(Q=P\) on (D). The relation
\((1-Q)/(1+Q)=l/F\) proves \eqref{answer} there.

\subsection*{All complex parameters and termination}
Lemma~\ref{convergence} makes the continued-fraction side meromorphic
on the connected domain
\(\{\Re x>0\}\times\mathbb C^3\). The gamma side is meromorphic there
as well. Applying the one-variable identity theorem successively
to the four variables extends equality from the real open box (D)
to this whole domain. This step also includes the removable cases
where \(D=0\) or an intermediate expression in the proof vanishes.

If \(n\) or \(\eta\) is an odd integer, then \(a_{|n|}=0\) or
\(a_{|\eta|}=0\), respectively. If \(l\) is a nonzero even integer,
then \(a_{|l|}=0\). In each case \(l/F\) is a finite rational
function of \(x\); its equality to the meromorphic gamma expression
on \(\Re x>0\) extends to all \(x\). Finally, \(l=0\) makes
\(P=1\) directly, and both sides of \eqref{answer} are identically
zero as meromorphic functions. This proves Theorem~\ref{main}.
\qed

\subsection*{Verification and attribution}
The attached verifier checks the two adjoint identities, the
transformed denominators, the gamma shift cancellation, and the
M\"obius identity symbolically. It also checks finite transformed
convergents exactly at rational parameter points and records
numerical consistency checks. The analytic convergence proof above,
not the finite checks, establishes the infinite assertion.

The Bauer--Muir method and gamma recurrence are classical inputs.
The source itself proves other formulas using that method, but
leaves the present formula as Conjecture 9.3. Our specific constant
adjoint factor and sign-changing shift supply the missing argument.

\section*{Task 2: sharp convergence rates}
Determine the leading error of the ordinary convergents for
\(\Re x>0\), uniformly on compact parameter sets away from poles.
Identify the parameter loci on which the leading term vanishes
and analyze the transition as \(\Re x\downarrow0\).

\section*{Task 3: certified accelerated evaluation}
Iterate the sign-changing Bauer--Muir shift to construct accelerated
approximants. Prove computable error bounds that remain useful near
the imaginary axis, and compare their cost with direct gamma evaluation.

\section*{Task 4: explicit continuants}
Find hypergeometric formulas for the numerator and denominator
continuants and identify their recurrence with an appropriate
polynomial family. Use those formulas to locate zeros in the
terminating cases and describe pole cancellation.

\section*{Task 5: a basic hypergeometric analogue}
Construct a \(q\)-continued fraction and a quotient of \(q\)-gamma
functions that tend to \eqref{answer} as \(q\uparrow1\).
Determine a convergence domain and an analogue of the sign-changing
shift \((x,\eta)\mapsto(x+2,-\eta)\).

\begin{thebibliography}{9}
\bibitem{CTZ}
X. Cao, Y. Tanigawa and W. Zhai,
\emph{Continued fraction formulae involving ratios of three gamma functions},
arXiv:2111.14143v1 (2021), 32 pp.
\url{https://arxiv.org/abs/2111.14143v1}.
The source of Conjecture 9.3 and the classical Bauer--Muir formulation
used above.
\end{thebibliography}
\end{document}
